\scanpage{24}
% L1–4
\par\noindent\begin{tabular}{|p{.92\linewidth}}
équivaut à la donnée d'une classe de tels $u$ mod addition d'un élément de $A$, i.e. à la donnée d'un élément de $C/A$ tel que, si $u$ en est un représentant, $(u)$ forme une $p$-base de
% L5–8
$C$ sur $A$\,; choisi $u$, la structure correspondante \deleted{\unsure{est}} est définie par la $A$-dérivation $D$ de $C$ telle que $Du=1$, et inversement, donné $D$, $u$ est \unsure{déterminé|déterminée} (mod $A$) par la condition $Du=1$.

% L9–12
Enfin, la donnée d'une \deleted{restriction des scalaires} \inserted{sous-$A$-hyperalgèbre} \inserted{$\unsure{U'|B'|\mathcal U'}$} de la $C$-hyperalgèbre tordue $\operatorname{Hom}_A(C,C)$, telle que \deleted{$\unsure{\operatorname{Hom}_A}$} $C$ devienne \deleted{\unsure{fibré}} homogène principal sous $U'$ et que $U'\simeq B$ \textendash\ équivaut à
% L13–16
la donnée d'une \deleted{\unsure{$x\in C$ $p$-base sur|$x\in C$ $p$-base de $u$}} \inserted{classe de} $p$-bases $(u)$ réduites à un élément $u$, en mettant dans la $=$ classe que $u$ les éléments $v=au+b$, avec $a\in A^*$ (unités de $A$) et $b\in A$\,; i.e.
% L17–20
l'ens de ces données est en correspondance biunivoque avec les éléments de l'espace \inserted{projectif} $\unsure{\mathbb P|P|\mathbf P}_A(C/A)=(C/A)/A^*$ associé au $A$-module $C/A$.
\end{tabular}
% archive sheet number: 23
